% GENERATED by tools/gen_appendices.py -- do not edit by hand.
% Re-run that script after adding, moving or retagging a chapter.
\chapter{Blueprint Coverage Matrix}\label{app:coverage}

Every topic in Cisco's published exam-topics document for CCNA 200-301
version~2.0, and the chapter of this volume that teaches it. The table is
generated from the \cmd{\textbackslash bp\{\}} tags in the chapter
sources and from the PDF itself, so it cannot drift from either.

\begin{longtable}{@{}L{1.3cm}L{9.9cm}L{3.3cm}@{}}
\toprule
\rowcolor{primary!15}
\textbf{Topic} & \textbf{Cisco's wording} & \textbf{Chapter} \\
\midrule
\endhead
\midrule
\multicolumn{3}{@{}l}{\textcolor{dom1!45!black}{\textbf{Domain 1.0 --- Network Infrastructure and Connectivity\quad 25\%}}} \\
\midrule
\tp{1.1} & Diagnose interface and cable (copper and fiber) issues such as collisions, errors, mismatched duplex, speed, distance, interface, signal levels, pin out, and cable types & 2, 3 \\
\rowcolor{lightbg}
\tp{1.2} & Describe the role and function of hypervisors, virtual machines, and containers & 29 \\
\tp{1.3} & Troubleshoot IPv4 address configuration, assignment, and subnetting (public and private) & 4 \\
\rowcolor{lightbg}
\tp{1.4} & Troubleshoot IPv6 address configuration, assignment, and prefix sizing (unicast and modified EUI 64) & 5 \\
\tp{1.5} & Describe wireless principles & 27 \\
\rowcolor{lightbg}
\tp{1.6} & Troubleshoot wired and wireless client connectivity (IP configuration, network reachability, and wireless security parameters on Windows, MacOS, and Linux) & 28 \\
\tp{1.7} & Troubleshoot DHCPv4 client, server, and relay on IOS devices & 20 \\
\midrule
\multicolumn{3}{@{}l}{\textcolor{dom2!45!black}{\textbf{Domain 2.0 --- Switching and Network Access\quad 25\%}}} \\
\midrule
\tp{2.1} & Configure network infrastructure connectivity (switch-to-switch and switch-to-router) & 10, 11 \\
\rowcolor{lightbg}
\tp{2.2} & Configure Layer 2 switch port attributes for edge-host connectivity (VLAN, PoE, port channel, and LACP) & 8, 9 \\
\tp{2.3} & Validate the accuracy of network documentation using CDP and LLDP & 13 \\
\rowcolor{lightbg}
\tp{2.4} & Troubleshoot basic Layer 2/Layer 3 connectivity and device operations using show commands (including show logs), ping, extended ping, trace route, and packet capture output & 6, 7, 14, 33 \\
\tp{2.5} & Configure operations of the Rapid Per VLAN Spanning Tree Protocol (Rapid PVST+) & 12 \\
\midrule
\multicolumn{3}{@{}l}{\textcolor{dom3!45!black}{\textbf{Domain 3.0 --- IP Routing\quad 20\%}}} \\
\midrule
\tp{3.1} & Interpret a routing table to identify the next hop for a packet (routing protocol, prefix/mask, administrative distance, metric, and default route) & 15 \\
\rowcolor{lightbg}
\tp{3.2} & Troubleshoot IPv4 and IPv6 static routing & 16 \\
\tp{3.3} & Configure single area OSPFv2 for IPv4 and OSPFv3 for IPv6 & 17, 18 \\
\rowcolor{lightbg}
\tp{3.4} & Interpret the operational status of First Hop Redundancy Protocols (HSRP and VRRP) & 19 \\
\midrule
\multicolumn{3}{@{}l}{\textcolor{dom4!45!black}{\textbf{Domain 4.0 --- Network Services and Security\quad 20\%}}} \\
\midrule
\tp{4.1} & Configure network devices with local usernames and as an AAA client (TACACS+ and RADIUS) for management & 23 \\
\rowcolor{lightbg}
\tp{4.2} & Manage device configuration and software files using secure file transfer operations with SFTP/SCP & 23 \\
\tp{4.3} & Configure NAT/PAT on IOS XE routers & 22 \\
\rowcolor{lightbg}
\tp{4.4} & Diagnose issues with DNS records (A, AAAA, CNAME, MX, NS, and PTR) to support host, web application, and mail server access by name & 21 \\
\tp{4.5} & Describe IPsec remote access and site-to-site VPNs (protocols and transport modes) & 26 \\
\rowcolor{lightbg}
\tp{4.6} & Configure IPv4 access control lists (standard, extended, numbered, and named ACLs) & 24 \\
\tp{4.7} & Configure Layer 2 security features & 25 \\
\midrule
\multicolumn{3}{@{}l}{\textcolor{dom5!45!black}{\textbf{Domain 5.0 --- AI, and Network Operations and Management\quad 10\%}}} \\
\midrule
\tp{5.1} & Describe the role of agentic AI in network operations & 32 \\
\rowcolor{lightbg}
\tp{5.2} & Select a prompt to send to a generative AI system to support network operations considering prompt components such as data classification, output format, persona, and instructions & 32 \\
\tp{5.3} & Describe network management approaches (device-based, cloud-based, controller- based, automation-based, and infrastructure as code) & 30 \\
\rowcolor{lightbg}
\tp{5.4} & Describe the function of SNMP in network operations & 30 \\
\tp{5.5} & Use configuration management mechanisms such as Ansible to execute commands & 31 \\
\rowcolor{lightbg}
\tp{5.6} & Interpret syslog message content, severity levels, and facilities & 30 \\
\bottomrule
\end{longtable}

